\chapter{Concurrency}\label{ch-par}

Algol 68 was one of the first languages to include concurrency in its
definition. A program can ask for several parts of itself to run at the
same time, and the parts can coordinate with \term{semaphores}, the
synchronization primitive that Edsger Dijkstra had introduced a few years
earlier. This chapter presents the two constructs, \code|par| and
\code|sema|, and uses them to solve three classic problems.

\begin{implnote}
Concurrency is the part of the language where the implementations differ
most. Algol 68 Genie runs the units of a \code|par| clause in threads of
the operating system. ga68 currently runs them one after the other, which
the Report allows, and does not yet support semaphores. a68j runs them in
threads when given the option \code|-fpar-threads|, and supports
semaphores only with that option; without it, a program that uses
\code|LEVEL| is rejected (``no lowering routine installed''). The
examples of this chapter were run with Algol 68 Genie and with a68j.
\end{implnote}

\section{The problems}\label{sec-par-problems}

Two problems arise as soon as parts of a program run at the same time and
share data. The first is \term{mutual exclusion}. If two parts both
execute \code|total +:= 1| on the same variable, each reads the old
value, adds one and stores the result; if both read before either stores,
one increment is lost. A sequence of operations that must not be
interleaved with another such sequence is a \term{critical section}, and
the program needs a way to let only one part at a time enter it.

The second problem is \term{communication}. One part, the producer,
computes values that another part, the consumer, needs. The consumer must
wait until a value is available, and if the values are passed through a
buffer of limited size, the producer must wait until there is room.

\section{Concurrent execution}\label{sec-par}

A list of expressions in parentheses, preceded by \code|par|, is a
\term{concurrent block} (\rr{parallel clause}). Its expressions may be
evaluated at the same time:

\example{ch08/par}

The \code|par| clause ends when all its expressions have ended; it has no
value. The expressions of an ordinary list may also be interleaved
(Section~\ref{sec-collateral}). What \code|par| adds is that each of its
expressions is a separate \term{process}: one of them can wait on a
semaphore while the others go on. Waiting outside a \code|par| clause has
no other process to wait for; the Report leaves it undefined, and Algol
68 Genie rejects \code|DOWN| and \code|UP| outside a \code|par| clause
altogether.

A \code|par| clause has a fixed number of expressions. A variable number
of processes is started by recursion, each level starting one worker and
the rest of the range:

\example{ch08/spawn}

\noindent Each worker writes only its own element of \code|squares|, so
the workers need no coordination, and the program prints
\code|1 4 9 16 25 36|.

\section{Semaphores}\label{sec-sema}

A semaphore is a counter that is changed only by two indivisible
operations. The type is \code|sema|, and a semaphore is created with the
operator \code|LEVEL|, which also gives its initial value:
\code|sema mutex = LEVEL 1|. The operations are unary operators:

\begin{itemize}
\item \code|DOWN s| waits until the counter of \code|s| is greater than
  zero and then decreases it by one. Dijkstra called this operation $P$;
  in today's libraries it is called \texttt{acquire} or \texttt{wait}.
\item \code|UP s| increases the counter by one, which may let a waiting
  part continue. Dijkstra called it $V$; today it is \texttt{release} or
  \texttt{signal}.
\end{itemize}

\code|LEVEL s| also gives the current value of the counter. A semaphore
whose counter is 0 or 1 is a \term{lock} (a mutex); a semaphore with a
larger counter limits the number of parts that may pass at the same
time.

The standard library declares \code|sema| as a record with a single field
that is a reference to an \code|int|, and \code|LEVEL n| allocates the
counter on the heap. A copy of a semaphore therefore shares the counter,
like the records with reference fields of Section~\ref{sec-ref-fields}:

\example{ch08/sema-copy}

\noindent prints 1. This is what makes it possible to pass a semaphore to
a procedure. A semaphore that was declared without \code|LEVEL|, such as
an element of \code|[5] sema forks|, holds no counter, and using it is a
run-time error; that is why the programs below initialize every
semaphore. \code|UP| lets all waiting parts compete for the counter again;
which one gets it is not specified, so a part may in principle wait
forever while others pass.

Note that \code|LEVEL|, like every unary operator, binds more tightly than
any binary operator: \code|LEVEL n - 1| means \code|(LEVEL n) - 1|, which
is not defined. The initial value \code|n - 1| must be written in
parentheses, as in the last program of this chapter.

\section{Mutual exclusion}\label{sec-mutex}

The following program lets three concurrent parts each add 100\,000 to a
shared total. Each increment is a critical section, protected by the
semaphore \code|mutex|:

\example{ch08/mutex}

Without the \code|DOWN| and \code|UP|, the result would usually be smaller
than 300\,000, and different from run to run, since increments would be
lost when two parts interleave. Such an error is called a \term{race
condition}; it is the classic problem of concurrent programming, and it
is hard to find precisely because it does not show up on every run.

The program takes the lock 300\,000 times, which costs time and makes the
parts wait for each other most of the time. Parts that share nothing need
no lock at all, so the better design lets each part count in a variable
of its own and takes the lock only once, to add its result:

\example{ch08/mutex-local}

A critical section should be as short as possible and should do no more
than the shared data require.

\section{Producer and consumer}\label{sec-buffer}

The second program passes values from a producer to a consumer through a
buffer with room for four values. Three semaphores coordinate the two
parts: \code|free| counts the free places, \code|full| the filled ones,
and \code|mutex| protects the buffer and its two positions.

\example{ch08/buffer}

\code|put| waits for a free place, stores the value and signals a filled
place; \code|get| does the opposite. If the producer is faster, it waits
in \code|DOWN free| when the buffer is full; if the consumer is faster,
it waits in \code|DOWN full| when the buffer is empty. Today's languages
offer the same structure ready-made as a blocking queue or a channel, but
it is built from exactly these parts.

The order of the two \code|DOWN|s in \code|put| matters. If the producer
took \code|mutex| first and then waited in \code|DOWN free| for a free
place, it would wait while holding the lock, and the consumer, which
needs the lock to make a place free, could never get it: a deadlock. The
rule is to wait for a resource before taking the lock that protects it,
and to hold the lock for as short a time as possible.

\section{The dining philosophers}\label{sec-philosophers}

Dijkstra's best-known synchronization problem has five philosophers
sitting at a round table, with one fork between each two of them. A
philosopher alternately thinks and eats, and needs both neighboring forks
to eat. If every philosopher picks up the left fork first and then waits
for the right one, all five may end up holding one fork each and waiting
forever, a \term{deadlock}.

One simple solution is a butler who lets at most four philosophers sit
at the table at the same time. Then at least one of them can always get
both forks. In the program, each fork is a semaphore with the initial
value 1, and the butler is a semaphore with the initial value 4:

\example{ch08/philosophers}

Each philosopher eats 1000 times and counts the meals in its own element
of \code|eaten|, which no other philosopher changes; the counts are
printed after the \code|par| clause has ended. Without the butler, the
program could hang.

\section*{Exercises}
\addcontentsline{toc}{section}{Exercises}

The \hyperref[ans-ch8]{answers} are at the end of the book.

\begin{enumerate}
\item What is the difference between \code|(a, b)| and \code|par (a, b)|?
\item What happens if a program executes \code|DOWN s| on a semaphore
  declared as \code|sema s = LEVEL 0|, and no other part ever executes
  \code|UP s|?
\item Is the semaphore \code|mutex| of the producer and consumer program
  necessary as long as there is exactly one producer and one consumer?
  When does it become necessary?
\item Give another solution of the dining philosophers that needs no
  butler. (Hint: let one philosopher pick up the forks in the other
  order.)
\item What happens to the producer and consumer program if the two
  \code|DOWN|s in \code|put| are exchanged, when the producer is faster
  than the consumer?
\item Why does \code|sema b = a; par (DOWN b, skip)| change
  \code|LEVEL a|?
\end{enumerate}

\section*{Programming problems}
\addcontentsline{toc}{section}{Programming problems}

\begin{enumerate}
\item Compute the sum of the elements of a large array with four
  concurrent parts, each summing a quarter.
\item Build a pipeline of three stages, each a concurrent part connected
  to the next by a buffer: the first generates numbers, the second
  squares them, the third prints them.
\item Implement a semaphore-protected counter as a record with a
  \code|sema| field and the procedures \code|increment| and \code|value|.
\item \textbf{A channel type.} Define a record type for a bounded buffer
  (a reference to an array, references to the two positions, and three
  semaphores) with procedures \code|send| and \code|receive|, so that
  several pipelines can be built from one definition.
\item \textbf{Readers and writers.} Several parts read a shared table
  and some change it. Readers may read at the same time, but a writer
  needs the table to itself. Solve it with semaphores, and think about
  whether a writer can be kept waiting forever.
\item \textbf{A barrier.} Build a barrier for $n$ processes from
  semaphores: no process continues past it until all have reached it.
  Make it reusable.
\item \textbf{Two transfers.} Two processes move money between the
  same two accounts in opposite directions, and each protects both
  accounts with a semaphore before it changes them. Show how they can
  deadlock, and avoid it by always taking the semaphore of the account
  with the smaller number first.
\item \textbf{Parallel merge sort.} Sort an array with merge sort in
  which the two halves are sorted in a \code|par| clause, down to a given
  depth, and compare the time with the sequential version.
\item \textbf{A prime sieve pipeline.} Generate the primes with a
  pipeline of processes: each process holds one prime, passes on the
  numbers it does not divide, and starts a new process for the first
  number that reaches it.
\end{enumerate}
